\documentclass[11pt,a4paper]{article}
\usepackage{fontspec,amsmath,amssymb,amsthm,microtype}
\setmainfont{lmroman10-regular.otf}[BoldFont=lmroman10-bold.otf,ItalicFont=lmroman10-italic.otf,BoldItalicFont=lmroman10-bolditalic.otf]
\usepackage[margin=25mm]{geometry}
\usepackage[unicode,hidelinks]{hyperref}
\hypersetup{pdfauthor={LangevinLean Project}}
\setlength{\emergencystretch}{3em}
\setlength{\parskip}{0.4em}
\setlength{\parindent}{0pt}
\allowdisplaybreaks[2]
\raggedbottom
\title{\texorpdfstring{The third-\allowbreak{}order one-\allowbreak{}step energy estimate for the actual Euler coupling}{The third-order one-step energy estimate for the actual Euler coupling}}
\author{LangevinLean Project}
\date{}
\begin{document}\maketitle
Draft mathematical exposition; these notes may contain mistakes.


The proof retains the cancellation in the defect cross term while bounding the complete defect by the existing first-order Brownian estimate. All gradient contraction and energy algebra concern the actual lifted gradient and actual integral solution.



In the \(L^2\) integral equation write \(g=\nabla f\), \(q=1-mh\), and \(e=\|X-Y_0\|_{L^2}\). The local defect and deterministic Euler difference give the exact decomposition


\[\Psi=(X-hg(X))-(Y_0-hg(Y_0)),\qquad D=\int_0^h\bigl(g(Y_t)-g(Y_0)\bigr)\,dt,\qquad X-hg(X)+B_h-Y_h=\Psi+D.\]


The proved pairing bound with constant \(R_{\rm gen}\), together with the score and Brownian moment bounds in dimension parameter \(d\), yields


\[\|X-hg(X)+B_h-Y_h\|_{L^2}^{2}\le q^2e^2+h^2qeR_{\rm gen}+\left(\frac32Lh\sqrt h\sqrt d\right)^2.\]

Earlier results used here: \cite[Section~3.13]{LangevinLean2026ULASC}.

\section{\texorpdfstring{Integrate the derived pairing bound and retain the exact cross coefficient}{Integrate the derived pairing bound and retain the exact cross coefficient}}


Let \ensuremath{\Psi} be the difference between the deterministic gradient steps of X and Y(0), and let D be the exact drift integral omitted by Euler. The integral equation gives the exact output error \ensuremath{\Psi}+D. On the admissible range \(hM\)\ensuremath{\le}1 and m\ensuremath{\le}M, the actual smooth strongly convex gradient step contracts by q=1\ensuremath{-}mh\ensuremath{\ge}0, so \ensuremath{\Vert}\ensuremath{\Psi}\ensuremath{\Vert}\ensuremath{\le}q\ensuremath{\Vert}X\ensuremath{-}Y(0)\ensuremath{\Vert}. The first-order total-defect theorem gives \ensuremath{\Vert}D\ensuremath{\Vert}\ensuremath{\le}(3/\allowbreak{}2)Mh\ensuremath{\surd}h\ensuremath{\surd}d. The separately proved weak-Dynkin pairing inequality bounds the pairing with the instantaneous gradient increment by t\ensuremath{\Vert}\ensuremath{\Psi}\ensuremath{\Vert}Rgen. Continuous linear integration then bounds \ensuremath{\langle\!\langle}\ensuremath{\Psi},D\ensuremath{\rangle\!\rangle} by h\ensuremath{{}^{2}}\ensuremath{\Vert}\ensuremath{\Psi}\ensuremath{\Vert}Rgen/\allowbreak{}2. Expanding \ensuremath{\Vert}\ensuremath{\Psi}+D\ensuremath{\Vert}\ensuremath{{}^{2}} and substituting the three bounds yields q\ensuremath{{}^{2}} times the original error squared, the exact cross term h\ensuremath{{}^{2}}q\ensuremath{\Vert}X\ensuremath{-}Y(0)\ensuremath{\Vert}Rgen, and the squared total-defect bound. The pairing premise is an intermediate analytic consequence of the actual stationary diffusion, discharged by the Dynkin and cutoff construction.



Earlier results used here: \cite[Section~3.11.1]{LangevinLean2026ULASC}; \cite[Section~3.11.3]{LangevinLean2026ULASC}; \cite[Section~3.13.2]{LangevinLean2026ULASC}; \cite[Section~3.13.5]{LangevinLean2026ULASC}; \cite[Section~3.13.6]{LangevinLean2026ULASC}.

\begin{thebibliography}{ULA-LSI}
\bibitem[ULA-SC]{LangevinLean2026ULASC}
LangevinLean Project. \emph{Unadjusted Langevin under Strong Convexity}. LL-01, 2026. Draft report, content revision 6f6f50decdb6. \url{https://odakyildiz.github.io/LangevinLean/docs/pdf/UnadjustedLangevin/FullProof.pdf}.
\end{thebibliography}

\end{document}
